\documentclass[11pt]{article}

\usepackage{amsmath, amssymb}
\usepackage{physics}
\usepackage{geometry}
\usepackage{bm}
\usepackage{enumitem}
\geometry{margin=1in}

\title{First Vertical Derivative (FVD) of Magnetic Data \\ 
with Non-Zero Observation Height}
\author{}
\date{}

\begin{document}

\maketitle

\section{Potential-Field Framework}

In a source-free region above magnetic sources, the magnetic field can be represented by a scalar potential $\Phi$:

\[
\mathbf{B} = -\nabla \Phi
\]

and satisfies Laplace's equation:

\[
\nabla^2 \Phi = 0.
\]

Expanding the Laplacian:

\[
\frac{\partial^2 \Phi}{\partial x^2}
+
\frac{\partial^2 \Phi}{\partial y^2}
+
\frac{\partial^2 \Phi}{\partial z^2}
=
0.
\]

\section{Fourier Transform in Horizontal Coordinates}

Taking the 2D Fourier transform in $x$ and $y$:

\[
\hat{\Phi}(k_x, k_y, z)
=
\mathcal{F}_{x,y}\{\Phi(x,y,z)\},
\]

the horizontal derivatives become:

\[
\frac{\partial^2}{\partial x^2} \rightarrow -k_x^2,
\qquad
\frac{\partial^2}{\partial y^2} \rightarrow -k_y^2.
\]

Thus Laplace's equation becomes

\[
\frac{\partial^2 \hat{\Phi}}{\partial z^2}
-
(k_x^2 + k_y^2)\hat{\Phi}
=
0.
\]

Define the radial wavenumber

\[
K = \sqrt{k_x^2 + k_y^2}.
\]

The governing equation reduces to

\[
\frac{\partial^2 \hat{\Phi}}{\partial z^2}
-
K^2 \hat{\Phi}
=
0.
\]

\section{Solution in the Vertical Direction}

The general solution is

\[
\hat{\Phi}(k_x,k_y,z)
=
A(k_x,k_y)e^{Kz}
+
B(k_x,k_y)e^{-Kz}.
\]

Requiring bounded behavior away from the sources eliminates the growing exponential term, yielding

\[
\hat{\Phi}(k_x,k_y,z)
=
B(k_x,k_y)e^{-Kz}.
\]

This exponential decay with height underpins upward continuation and derivative operators.

\section{Vertical Derivative Operator}

Let $B(x,y,z)$ be a magnetic component measured at elevation $z$.
In Fourier space,

\[
\hat{B}(k_x,k_y,z)
=
\hat{B}(k_x,k_y,0)\,e^{-Kz}.
\]

Taking the vertical derivative,

\[
\frac{\partial \hat{B}}{\partial z}
=
- K \hat{B}(k_x,k_y,0)e^{-Kz}.
\]

Evaluated at the observation plane $z = z_{\text{obs}}$:

\[
\boxed{
\mathcal{F}
\left\{
\left.
\frac{\partial B}{\partial z}
\right|_{z=z_{\text{obs}}}
\right\}
=
- K \hat{B}(k_x,k_y,z_{\text{obs}})
}
\]

Thus, at any constant elevation plane,

\[
\boxed{
\frac{\partial B}{\partial z}
=
\mathcal{F}^{-1}
\left\{
- K \hat{B}
\right\}.
}
\]

In practice, the sign depends on the $z$-axis convention and is often absorbed into implementation.

\section{Upward Continuation}

Upward continuation by $\Delta z > 0$ is given by

\[
\hat{B}(z+\Delta z)
=
\hat{B}(z)e^{-K\Delta z}.
\]

\section{Stabilized First Vertical Derivative}

Because the operator multiplies by $K$, high-frequency noise is amplified.
A stabilized form is obtained by first upward continuing:

\[
\boxed{
\widehat{\text{FVD}}(z+\Delta z)
=
- K e^{-K\Delta z} \hat{B}(z).
}
\]

This corresponds physically to computing the vertical derivative at a higher datum.

\section{Handling Variable Observation Height}

If data are collected on a variable topographic or draped surface:

\begin{enumerate}[leftmargin=1.5em]
\item Continue data to a constant elevation plane (preferably upward for stability).
\item Apply the FVD operator in Fourier space.
\end{enumerate}

An alternative is to compute an equivalent source representation and evaluate the field on a chosen datum prior to differentiation.

\section{One-Dimensional Case}

For a profile $B(x)$,

\[
K = |k|,
\]

and

\[
\boxed{
\mathcal{F}
\left\{
\frac{\partial B}{\partial z}
\right\}
=
- |k| \hat{B}(k).
}
\]

\end{document}